\section*{Bab \ref{ch:g7:geological-time} --- Umur Bumi}

\begin{solution}{exo:g7:geological-time:1}
(a) \omterm{def:g7:geological-time:unconformity}{ketidakselarasan}; (b) inklusi; (c) superposisi; (d) suksesi \omterm{def:g7:sediments-and-fossils:fossil}{fosil};
(e) hubungan potong-memotong.
\end{solution}

\begin{solution}{exo:g7:geological-time:2}
(a) Batu pasir merah, di bagian bawah, adalah yang paling tua; batu pasir
kuning, di bagian atas, adalah yang paling muda. (b) \omterm{def:g7:geological-time:superposition}{Prinsip superposisi}.
(c) Lapisan-lapisan itu belum terbalik sejak terbentuk.
\end{solution}

\begin{solution}{exo:g7:geological-time:3}
(a) Paleozoikum, Mesozoikum, Kenozoikum. (b) Kenozoikum. (c) Mesozoikum.
(d) Jura datang tepat sesudah Trias; Kapur datang tepat sebelum Paleogen.
\end{solution}

\begin{solution}{exo:g7:geological-time:4}
(a) Salah: \omterm{def:g7:geological-time:relative-dating}{penanggalan relatif} hanya mengurutkan \omterm{def:g3:stones-pebbles-sand:rock}{batuan}; umur dalam tahun
adalah \omterm{def:g7:geological-time:relative-dating}{umur absolut}. (b) Benar. (c) Salah: \omterm{def:g3:stones-pebbles-sand:pebble}{kerakal-kerakal} itu lebih tua
daripada konglomerat yang mengandungnya (\omterm{def:g7:geological-time:inclusion}{prinsip inklusi}). (d) Benar
(\omterm{def:g7:geological-time:fossil-succession}{prinsip suksesi fosil}).
\end{solution}

\begin{solution}{exo:g7:geological-time:5}
(a) Lebih muda: granit itu memotong lapisan 3, dan \omterm{def:g3:stones-pebbles-sand:rock}{batuan} yang memotong
\omterm{def:g3:stones-pebbles-sand:rock}{batuan} lain lebih muda (potong-memotong). (b) Lebih tua: lapisan 4 terletak
di atasnya tanpa terpotong. (c) Lapisan 1, lapisan 2, lapisan 3, granit,
lapisan 4. (d) \omterm{def:g3:stones-pebbles-sand:rock}{Batuan} di atasnya terkikis oleh \omterm{def:g5:water-shapes-the-land:erosion}{erosi} sampai granit itu
mencapai permukaan; lapisan 4 kemudian diendapkan di atas permukaan \omterm{def:g5:water-shapes-the-land:erosion}{erosi}
ini, sebuah \omterm{def:g7:geological-time:unconformity}{ketidakselarasan}.
\end{solution}

\begin{solution}{exo:g7:geological-time:6}
(a) $538.8 - 251.9 = 286.9$ juta tahun. (b) $251.9 - 66 = 185.9$
juta tahun. (c) 66 juta tahun. (d) Paleozoikum adalah yang paling panjang,
Kenozoikum yang paling pendek (sampai sekarang).
\end{solution}

\begin{solution}{exo:g7:geological-time:7}
(a) $4540 - 538.8 = 4001.2$ juta tahun. (b) $4001.2 \div 4540
\approx 0.881$: kira-kira \textbf{88~\%}. (c) Sebelum Paleozoikum, hewan
bercangkang belum umum, dan makhluk hidup yang lunak hampir tidak pernah
menjadi \omterm{def:g7:sediments-and-fossils:fossil}{fosil}; cangkang dan tulang, yang mudah menjadi \omterm{def:g7:sediments-and-fossils:fossil}{fosil}, baru umum
sejak Paleozoikum.
\end{solution}

\begin{solution}{exo:g7:geological-time:8}
(a) $4540 \div 100 = 45.4$ juta tahun per meter. (b) $538.8 \div 45.4
\approx 11.9$: kira-kira \qty{11.9}{m} sebelum garis akhir. (c) $66 \div
45.4 \approx 1.45$: kira-kira \qty{1.45}{m}. (d) $0.3 \div 45.4 \approx
0.0066$~m, yaitu kira-kira \qty{0.66}{cm}, kurang dari satu sentimeter:
seorang pelari hampir tidak akan memperhatikannya.
\end{solution}

\begin{solution}{exo:g7:geological-time:9}
(a) Silur lebih dahulu. Ya: \omterm{def:g3:stones-pebbles-sand:rock}{batuan} yang lebih tua terletak di bawah \omterm{def:g3:stones-pebbles-sand:rock}{batuan}
yang lebih muda. (b) Lapisan-lapisan kelabu diendapkan secara mendatar;
lapisan-lapisan itu didorong sampai tegak; lapisan-lapisan itu terkikis;
\omterm{def:g3:stones-pebbles-sand:sand}{pasir} merah diendapkan di atasnya. (c) Permukaan itu adalah permukaan \omterm{def:g5:water-shapes-the-land:erosion}{erosi}
di antara \omterm{def:g3:stones-pebbles-sand:rock}{batuan-batuan} yang umurnya sangat berbeda: waktu yang lama ketika
lapisan-lapisan kelabu itu diangkat dan terkikis tidak meninggalkan satu
lapisan pun.
\end{solution}

\begin{solution}{exo:g7:geological-time:10}
(a) Kedua serpih dengan amonit berumur sama, demikian pula kedua batu kapur
dengan bulu babi: \omterm{def:g7:sediments-and-fossils:fossil}{fosil} yang sama, umur yang sama (\omterm{def:g7:geological-time:fossil-succession}{prinsip suksesi fosil}).
(b) Yang paling tua adalah batu pasir dengan trilobit di tebing A (batu
pasir itu terletak di bawah serpih dengan amonit, yang sama tuanya dengan
lapisan terbawah tebing B). Yang paling muda adalah lempung dengan daun di
tebing B. (c) Dari bawah ke atas: batu pasir dengan trilobit; serpih dengan
amonit; batu kapur dengan bulu babi; lempung dengan daun.
\end{solution}

\begin{solution}{exo:g7:geological-time:11}
(a) Sekis. Sekis terletak di bawah konglomerat (superposisi), dan \omterm{def:g3:stones-pebbles-sand:pebble}{kerakal}
konglomerat itu berupa kepingan sekis (inklusi); permukaan \omterm{def:g5:water-shapes-the-land:erosion}{erosi} juga
memotong sekis. (b) Sesudahnya: batu pasir itu diendapkan secara mendatar,
jadi batu pasir itu dimiringkan kemudian. (c) Sebelumnya: lempung itu
menutupi \omterm{def:g7:earthquakes:fault}{sesar} tanpa terpatahkan. (d) Sekis terbentuk; sekis terkikis
(permukaan yang bergelombang); konglomerat, batu pasir, dan batu kapur
diendapkan, dengan urutan itu; lapisan-lapisan itu dimiringkan, dan \omterm{def:g7:earthquakes:fault}{sesar}
mematahkannya --- penampang itu sendiri tidak mengatakan mana dari kedua
peristiwa ini yang lebih dahulu; lapisan-lapisan yang miring dan \omterm{def:g7:earthquakes:fault}{sesar} itu
terkikis rata; lempung diendapkan di atasnya. (e) Dua: permukaan yang
bergelombang di antara sekis dan konglomerat, dan permukaan di antara
lapisan-lapisan yang miring dan lempung yang mendatar.
\end{solution}

\begin{solution}{exo:g7:geological-time:12}
(a) Separuh dari 365 hari adalah 182.5 hari. Januari sampai Juni berjumlah
$31+28+31+30+31+30 = 181$ hari, jadi sejarah itu sudah lewat separuhnya
1.5 hari memasuki bulan Juli: pada \textbf{2~Juli pukul 12.00 siang}.
(b) $4540 \div 2 = 2270$ juta tahun. (c) $80 \div \num{4540000000} \times \num{31536000} \approx 0.56$:
kira-kira \textbf{0.6~detik}.
\end{solution}

\begin{solution}{pb:g7:geological-time:1}
\textbf{1.} \num{4540} juta tahun.

\textbf{2.} $4540 \div 365 \approx 12.4$ juta tahun.

\textbf{3.} $\num{4540000000} \div (365 \times 24) = \num{4540000000}
\div 8760 \approx \num{518000}$ tahun (dengan nilai bulat \num{12.4}
juta, $\num{12400000} \div 24 \approx \num{517000}$).

\textbf{4.} $\num{4540000000} \div \num{525600} \approx 8638$: kira-kira
\num{8600} tahun.

\textbf{5.} $4540 - 4030 = 510$ juta tahun sesudah Bumi terbentuk;
$510 \div 12.4 \approx 41$ hari. Januari mempunyai 31 hari, jadi 41 hari
sudah berlalu pada awal tanggal \textbf{11~Februari}.

\textbf{6.} $4540 - 3500 = 1040$ juta tahun; $1040 \div 12.4 \approx
83.9$ hari (83.6 dengan nilai yang tidak dibulatkan, $12.44$). Januari dan
Februari berjumlah 59 hari: 24.6 sampai 24.9 hari memasuki bulan Maret,
pada tanggal \textbf{25~Maret}.

\textbf{7.} $66 \div 12.4 \approx 5.3$ hari sebelum akhir tahun, yaitu
$365 - 5.3 = 359.7$ hari sesudah 1~Januari. Sebelas bulan pertama
berjumlah 334 hari: 25.7 hari memasuki bulan Desember, pada tanggal
\textbf{26~Desember}.

\textbf{8.} $538.8 \div 4540 \approx 0.119$: \textbf{11.9~\%} dari tahun
itu yang tersisa, yaitu $0.119 \times 365 \approx 43$ hari. $365 - 43 =
322$ hari sesudah 1~Januari: Paleozoikum dimulai pada bulan
\textbf{November} (pada tanggal 18~November).

\textbf{9.} $\num{10000} \div 8600 \approx 1.2$ menit.

\textbf{10.} $5000 \div 8600 \approx 0.58$ menit, dan $0.58 \times 60
\approx 35$: kira-kira \textbf{35 detik} sebelum tengah malam.

\textbf{11.} $\num{300000} \div 8600 \approx 34.9$: kira-kira \textbf{35
menit} sebelum tengah malam.

\textbf{12.} 35 menit sebelum tengah malam tanggal 31~Desember: spesies kita
muncul kira-kira \textbf{pukul 23.25}, pada jam terakhir tahun itu.
\end{solution}
